% Generated draft: requires mathematical and source review.
\documentclass[11pt]{article}
\usepackage[margin=25mm]{geometry}
\usepackage{fontspec}
\setmainfont{Latin Modern Roman}
\usepackage{amsmath,amssymb,mathtools,unicode-math}
\setmathfont{Latin Modern Math}
\usepackage{enumitem,needspace,xurl,hyperref}
\setlength{\parindent}{0pt}
\setlength{\parskip}{5pt}
\setlength{\emergencystretch}{3em}
\newcommand{\N}{\mathbb{N}}
\newcommand{\Z}{\mathbb{Z}}
\newcommand{\Q}{\mathbb{Q}}
\newcommand{\R}{\mathbb{R}}
\newcommand{\C}{\mathbb{C}}
\newcommand{\F}{\mathbb{F}}
\newcommand{\A}{\mathbb{A}}
\newcommand{\PP}{\mathbb{P}}
\newcommand{\E}{\mathbb{E}}
\newcommand{\eps}{\varepsilon}
\newcommand{\dd}{\,\mathrm{d}}
\begin{document}
\section*{UnsolvedMath ID: 30000517}
\textbf{Source problem code:} catalogue entry
\subsection{Definitions and mathematical statement}
Let $\Sigma\subset\R^3$ be the intended shell midsurface, let $x_0\in\Sigma$ be the load point, and let $h>0$ denote thickness, tending to zero relative to a fixed geometric scale. If $n(x_0)$ is a unit normal and $F$ is the load amplitude, a normal point force is represented by $F n(x_0)\delta_{x_0}$, acting on a smooth vector test field $\varphi$ by $F n(x_0)\cdot\varphi(x_0)$. Let $u_h$ denote the displacement solution in the shell model intended by the source; its normal component is $u_h\cdot n$.

The problem asks to determine the small-thickness asymptotic behavior of this solution near $x_0$, with particular attention to hyperbolic and parabolic geometries. In the usual differential-geometric terminology, a hyperbolic point has Gaussian curvature $K(x_0)<0$, whereas a parabolic point has $K(x_0)=0$ and rank-one second fundamental form. This distinguishes a parabolic point from a planar point, where the second fundamental form has rank zero.

The requested local asymptotics should identify the relevant spatial scales, leading displacement behavior and its range of validity as $h\to0$. The source supplies no shell equations, constitutive law, boundary conditions, load scaling with $h$, or solution concept. Accordingly, no particular operator equation, power of $h$, isotropic rescaling, leading profile or remainder estimate can be asserted from the extract. The coupled limit of approaching the point load and decreasing thickness, and the intended curvature nondegeneracy assumptions in the parabolic case, also require specification.
\subsection{Short English statement}
Find the displacement asymptotics close to a normal point force as a shell becomes thin, particularly at hyperbolic and parabolic points. The governing shell model and limit conventions are missing, so the draft retains the local asymptotic target without inventing an equation or scaling law.
\subsection{Sources}
\begin{itemize}
\item Supplied report: \emph{Applications of Asymptotic Analysis} (2007), Oberwolfach Reports, DOI 10.4171/owr/2006/28, pp. 1663--1730; proposed year 2006. \url{https://doi.org/10.4171/owr/2006/28}.
\item Supplied reference labelled ``shell asymptotic study'': \url{https://www.sciencedirect.com/science/article/pii/S0045794902000093}.
\item Local UnsolvedMath registry, permanent ID 30000517, catalogue entry; created 25 August 2026, updated 16 September 2026. It labels the statement exact and reports comparison with the official report. Its 21 August 2026 literature triage labels the question open: selected shallow-shell geometries have asymptotic and numerical studies, while the general local point-load problem, especially hyperbolic and parabolic cases, is reported unsettled. This dated assessment requires independent verification.
\end{itemize}
\end{document}
